% ============================================================================
% tex/frontmatter_content/abstract.tex
%
% Content-only file used by \AbstractPage{...} in main.tex.
%
% House rule:
%   - Do NOT add \chapter/\section here.
%   - Do NOT include \begin{document}.
%   - Keep this file as pure content that can be injected into a macro argument.
% ============================================================================

	One of the main tasks for computer hardware and software is \textbf{lexical input} -- sending
	\underline{elements} of human language \underline{input sets} to computer applications like word processors,
	spreadsheets, etc. Input \underline{elements} include characters (letters, numbers, symbols, etc.) and
	(possibly single-word) phrases comprised of characters. \underline{Input sets} include languages like
	English or Chinese and specialized input sets like the International Phonetic Alphabet (IPA).

	The standard QWERTY keyboard may not best serve as lexical input system in a variety of situations:

	\begin{itemize}
		\item It may be absent, as for tablet computers.
		\item It may be physically difficult to use, as for rheumatoid arthritis sufferers.
		\item Its fixed markings may not suit input sets, as for several Far Eastern languages.
	\end{itemize}

	This research describes a device-agnostic, alternative input system (µLex™) that may better serve in these
	situations alongside, or in place of, keyboards:

	\begin{itemize}
		\item Keyboard absence and the limitations inherent in fixed markings are countered by a display-resident,
		variable-contents presentation/selection grid.
		\item Physical difficulties, inherent in mouse-clicking as well as keyboarding, is countered by alternative
		actuation methods.
		\item Cursor-based slower operation, inherent in on-display input systems, is countered by presenting
		user-friendly, shorter selection sequences that result in longer length input.
	\end{itemize}

	Design concepts and sample apps based on the design concepts are described to illustrate how they are combined
	to increase efficiency for representative, increasingly complex input sets (Spanish letters and punctuation
	marks, Pan-Euro Characters, IPA, Chinese).

	The thesis for this research is that gesture efficient input methods for global languages can be generated
	algorithmically from phrase and character frequency lists. The (language-independent) measure for efficiency
	is Gestures per Lexeme (lexeme $\sim$ word or phrase) and proof is by construction of the algorithm.

	It is shown that the apps are efficient; that is, they optimize the number of input characters for the number
	of user gestures.

	Further it is shown that the efficiency measures hold regardless of the peripheral devices (e.g., keyboards,
	mice, touchscreens, etc.) used for input.

	The benefits of increasing efficiency are discussed with respect to assistive technologies, non-English input,
	and novel, special-purpose peripherals.

	An app typology is proposed to assist in determining how to design apps for new input sets.
